Tenim dues ones monocromàtiques, una de $\lambda = 750$ nm, corresponent al color vermell,
i una altra de $\lambda = 550$ nm, corresponent al color verd.

\begin{apartat}{1,25}
Quina de les dues ones té els fotons més energètics? Calculeu la freqüència i l'energia
d'un fotó i les d'un mol de fotons per a cadascuna de les dues radiacions.
\end{apartat}

\begin{apartat}{1,25}
A continuació, volem reproduir l'efecte fotoelèctric il·luminant amb llum monocromàtica
una placa de rubidi. Determineu la longitud d'ona llindar perquè es produeixi l'efecte
fotoelèctric. Les ones visibles de l'apartat anterior seran capaces d'arrencar un electró
de la superfície del rubidi? Si això és possible, quina serà l'energia cinètica adquirida
per l'electró?

Per a l'efecte fotoelèctric, representeu esquemàticament com varia l'energia cinètica
màxima dels electrons arrencats en funció de l'energia dels fotons incidents. Comenteu el
significat de les dues zones diferenciades.
\end{apartat}

\dades{%
  La funció de treball del rubidi és $3,46\times 10^{-19}$ J.\\
  $N_A = 6,022\times 10^{23}$.\\
  $c = 3,00\times 10^{8}\ \mathrm{m\,s^{-1}}$.\\
  $h = 6,63\times 10^{-34}\ \mathrm{J\,s}$.\\
  $m_\mathrm{e} = 9,11\times 10^{-31}\ \mathrm{kg}$.}
